\documentclass[12pt]{article}

\usepackage[margin=1in]{geometry}
\usepackage[T1]{fontenc}
\usepackage{lmodern}
\usepackage{microtype}
\usepackage{amsmath,amssymb,mathtools}
\usepackage{booktabs}
\usepackage{setspace}
\usepackage{hyperref}
\usepackage{xurl}

\hypersetup{
  colorlinks=true,
  linkcolor=black,
  citecolor=black,
  urlcolor=blue
}

\onehalfspacing
\widowpenalty=10000
\clubpenalty=10000

\title{The Gradient of the Inverse Clock Rate:\\
Static Acceleration, the Weak Field, and Schwarzschild Spacetime}
\author{Light-Meter Physics}
\date{}

\begin{document}

\maketitle

\begin{abstract}
For observers held at rest in a static spacetime, the spatial gradient of the logarithm of the lapse gives their proper acceleration. This paper expresses the relation using the dimensionless inverse clock-rate ratio $\delta\tau\equiv\mathrm{d}t/\mathrm{d}\tau$ and the spatial gradient $\nabla$. With the gravitational direction defined inward, $\mathbf{g}=c^2\nabla\ln\delta\tau$, while the acceleration required to hold position is $\mathbf{a}_{\mathrm{hold}}=-\mathbf{g}$. In the weak-field limit this becomes Newtonian gravity, and in the Schwarzschild exterior it gives the exact local acceleration measured by a supported observer. The paper distinguishes this static-field identity from a weak-field rate ledger for a moving clock, in which both gravitational potential and kinetic energy contribute. An idealized Sun--Earth example illustrates the reference and model dependence of that ledger. Clock-comparison experiments, the horizon limit, and acceleration-associated scales are discussed with explicit qualifications. The term ``temporal density'' is treated only as project shorthand for an inverse clock-rate ratio, not as a new physical field or a material density.
\end{abstract}

\noindent\textbf{Keywords:} gravitational time dilation, lapse, clock rates, static observers, proper acceleration

\section{Introduction}

A clock held at a fixed position outside a spherical gravitating body accumulates less proper time per unit reference time than a clock at infinity. In Schwarzschild spacetime this comparison is described by the lapse $\alpha(r)=\sqrt{1-r_s/r}$, where $r_s=2GM/c^2$. The reciprocal rate is useful for expressing the same comparison in the opposite direction. In this paper it is denoted
\begin{equation}
  \delta\tau\equiv\frac{\mathrm{d}t}{\mathrm{d}\tau},
  \qquad
  \epsilon_{\mathrm{grav}}\equiv\delta\tau-1\quad\text{(static case)}.
  \label{eq:inverse-rate}
\end{equation}
For the static Schwarzschild comparison, $t$ is normalized to proper time at infinity, so $\delta\tau=1/\alpha$, and $\tau$ is the proper time of the held clock. Despite its symbol, $\delta\tau$ is a dimensionless rate ratio, not an elapsed proper-time interval. The informal label ``temporal density'' is used here only as a mnemonic for this inverse rate; it does not denote a new matter density or independent field.

The subject here is the spatial change of that rate ratio. In a static spacetime, the gradient of $\ln\delta\tau$ is related exactly to the acceleration required to remain at rest. Its magnitude can be expressed as an acceleration by multiplying by $c^2$. With the gravitational direction taken inward, the relation has the compact form $\mathbf{g}=c^2\nabla\ln\delta\tau$. This is a standard kinematic identity for a specified family of static observers, not an additional law of gravity \cite{Wald,TaylorWheeler,SantiagoVisser}.

The distinction between a clock comparison and an accelerometer reading is essential. A clock-rate ratio compares worldlines relative to a chosen time coordinate. An accelerometer carried by a supported observer records that observer's proper acceleration. The gradient connects these concepts for static observers, but does not imply that an individual clock locally feels a rate offset. A freely falling observer follows a geodesic and reads zero proper acceleration.

The analysis first defines the rate ratio and spatial gradient, then derives the acceleration identity, and recovers Newtonian gravity in the weak-field limit. The exact acceleration identity applies to static observer families in a general static spacetime; its explicit radial formulas below specialize to the Schwarzschild vacuum exterior of a spherical, nonrotating source. A separate moving-clock ledger, terrestrial clock comparisons, and the limiting behavior near a horizon follow. Rotating or time-dependent spacetimes and arbitrary moving observers require additional structure.

\section{Static observers, reference time, and notation}

A general static spacetime can locally be written in block-diagonal form
\begin{equation}
  \mathrm{d}s^2=-\alpha(\mathbf{x})^2c^2\mathrm{d}t^2
    +h_{ij}(\mathbf{x})\,\mathrm{d}x^i\mathrm{d}x^j,
  \label{eq:static-metric}
\end{equation}
where $\alpha>0$ is the lapse and $h_{ij}$ is the spatial metric on a constant-$t$ slice. A static observer has fixed spatial coordinates, so the interval along that worldline gives
\begin{equation}
  \mathrm{d}\tau=\alpha\,\mathrm{d}t,
  \qquad
  \delta\tau=\alpha^{-1}.
  \label{eq:static-rate}
\end{equation}
For an asymptotically flat spacetime, the normalization $\alpha\to1$ at infinity makes $t$ agree there with the proper time of a stationary observer. This choice gives $\delta\tau$ a definite reference. In a more general static spacetime without such a boundary, the normalization must instead be specified by a reference observer or convention.

Throughout this paper, $\nabla f$ denotes the spatial gradient of a scalar $f$ formed with the spatial geometry $h_{ij}$ and expressed in the local orthonormal frame of static observers. It therefore measures change per unit proper spatial distance, not per unit coordinate distance. In Schwarzschild coordinates, for a radial scalar,
\begin{equation}
  \nabla f=\alpha\,\frac{\mathrm{d}f}{\mathrm{d}r}\,\hat{\mathbf{r}},
  \label{eq:radial-gradient}
\end{equation}
where $\hat{\mathbf{r}}$ points outward. This convention keeps the gradient in the paper tied to a physical ruler.

The logarithm is useful for two related reasons. First, it converts multiplicative rate comparisons into additive differences:
\begin{equation}
  \ln\!\left(\frac{\delta\tau_2}{\delta\tau_1}\right)
  =\ln\delta\tau_2-\ln\delta\tau_1.
  \label{eq:log-comparison}
\end{equation}
Second, a constant rescaling of the reference time, $t'=k t$, multiplies $\delta\tau$ by $k$ and adds the constant $\ln k$ to $\ln\delta\tau$. Its spatial gradient is unchanged. Thus the local gradient does not depend on the arbitrary constant normalization of the static time coordinate, although the rate ratio itself does depend on the chosen reference.

A field such as $\nabla\ln\delta\tau$ describes variation across a specified family of static observers. It is not inferred from a single absolute clock reading. Locally, a supported observer can measure its own proper acceleration with an accelerometer; comparing static clocks at neighboring locations measures the corresponding rate gradient.

\section{From the inverse-rate gradient to acceleration}

The static observers follow the timelike Killing flow associated with $t$. Their four-velocity is the normalized time-translation direction. The covariant four-acceleration of this flow is the spatial gradient of the logarithm of the lapse \cite{Wald,SantiagoVisser}. Restoring ordinary acceleration units gives the proper acceleration vector required to hold a static position:
\begin{equation}
  \mathbf{a}_{\mathrm{hold}}
  =c^2\nabla\ln\alpha.
  \label{eq:hold-alpha}
\end{equation}
Since $\delta\tau=1/\alpha$, Eq.~\eqref{eq:hold-alpha} is equivalently
\begin{equation}
  \boxed{\mathbf{a}_{\mathrm{hold}}
  =-c^2\nabla\ln\delta\tau.}
  \label{eq:hold-delta}
\end{equation}
The sign identifies the physical direction. For an ordinary attractive spherical source, $\alpha$ increases outward and $\delta\tau$ decreases outward. The support acceleration therefore points outward, away from the source.

It is useful to define the inward gravitational direction for the static observer family as the negative of the acceleration needed to hold position:
\begin{equation}
  \boxed{\mathbf{g}\equiv-\mathbf{a}_{\mathrm{hold}}
  =c^2\nabla\ln\delta\tau.}
  \label{eq:g-definition}
\end{equation}
The magnitude of either vector is the local hovering acceleration,
\begin{equation}
  g=a_{\mathrm{hold}}=c^2\left|\nabla\ln\delta\tau\right|.
  \label{eq:g-magnitude}
\end{equation}
The vector in Eq.~\eqref{eq:g-definition} is defined relative to the specified static observer family; in a gravitational setting it gives the inward direction of free fall in that frame. A physical accelerometer attached to a supported observer points oppositely, recording the support acceleration in Eq.~\eqref{eq:hold-delta}. It is not a universal formula for the coordinate acceleration of an arbitrary moving observer.

Equations~\eqref{eq:hold-delta}--\eqref{eq:g-magnitude} apply to static observers in static spacetime. They do not say that every observer's motion is determined by a clock-rate gradient, nor that the lapse is itself a force. For a freely falling observer the four-acceleration vanishes, even while the observer's clock rate relative to the static coordinate time may differ from that of a static clock.

\subsection{Flat spacetime and observer acceleration}

Rindler coordinates describe a Born-rigid family of accelerated observers in flat spacetime; its members do not all have the same proper acceleration. In one Rindler wedge, the metric and lapse can be written \cite{Rindler}
\begin{equation}
  \mathrm{d}s^2=-\alpha(z)^2c^2\mathrm{d}t^2
    +\mathrm{d}x^2+\mathrm{d}y^2+\mathrm{d}z^2,
  \qquad
  \alpha(z)=1+\frac{a_0z}{c^2},
  \quad z>-\frac{c^2}{a_0},
  \label{eq:rindler-metric}
\end{equation}
where $a_0>0$ is the proper acceleration of the observer at $z=0$. The static observers have
\begin{equation}
  \mathbf{a}_{\mathrm{hold}}
  =c^2\nabla\ln\alpha
  =\frac{a_0}{1+a_0z/c^2}\,\hat{\mathbf{z}}.
  \label{eq:rindler-acceleration}
\end{equation}
The Riemann curvature tensor nevertheless vanishes: the nonzero lapse gradient describes this accelerated observer family, not spacetime curvature. Thus a lapse gradient alone does not establish a new gravitational mechanism or a material field.

\section{The weak-field limit and Newtonian gravity}

Let $\Phi$ be the Newtonian gravitational potential, with $\Phi(\infty)=0$ and $\Phi<0$ near an attractive source. In a weak, static field,
\begin{equation}
  \alpha\simeq1+\frac{\Phi}{c^2},
  \qquad
  \delta\tau\simeq1-\frac{\Phi}{c^2},
  \qquad
  \epsilon_{\mathrm{grav}}\simeq-\frac{\Phi}{c^2}.
  \label{eq:weak-rate}
\end{equation}
Taking the spatial gradient and retaining the leading term gives
\begin{equation}
  \mathbf{g}=c^2\nabla\ln\delta\tau
  \simeq-\nabla\Phi,
  \label{eq:newtonian-limit}
\end{equation}
which is the Newtonian gravitational field. The logarithm of the inverse rate is dimensionless; its gradient has units of inverse length, and multiplication by $c^2$ gives acceleration.

Near Earth's surface, with $g\simeq9.81\ \mathrm{m\,s^{-2}}$, the magnitude of the fractional rate gradient is approximately
\begin{equation}
  \left|\nabla\ln\delta\tau\right|
  \simeq\frac{g}{c^2}
  \simeq1.09\times10^{-16}\ \mathrm{m^{-1}}.
  \label{eq:earth-gradient}
\end{equation}
Over a small upward height $h$, the upper static clock runs faster, with
\begin{equation}
  \frac{\Delta f}{f}\simeq\frac{gh}{c^2},
  \qquad
  \Delta\ln\delta\tau\simeq-\frac{gh}{c^2}.
  \label{eq:height-redshift}
\end{equation}
For $h=1\ \mathrm{m}$ the fractional rate difference is about $1.1\times10^{-16}$. Optical-clock comparisons have tested gravitational redshift over height differences of tens of centimeters, and strontium-clock measurements have resolved the predicted frequency gradient within a millimeter-scale atomic sample \cite{Chou,Bothwell}. These are relative frequency comparisons, not measurements of an absolute local clock-rate offset.

\section{A separate ledger for moving clocks}

The static identity above should not be confused with the weak-field rate of a moving clock. For this calculation, use $t$ as the time coordinate of a specified weak-field coordinate system with its stated asymptotic normalization. The leading post-Newtonian relation for a slowly moving clock is
\begin{equation}
  \frac{\mathrm{d}\tau}{\mathrm{d}t}
  \simeq1+\frac{\Phi}{c^2}-\frac{v^2}{2c^2},
  \qquad
  \delta\tau-1
  \simeq-\frac{\Phi}{c^2}+\frac{v^2}{2c^2},
  \label{eq:moving-ledger}
\end{equation}
where $\Phi$ is the total Newtonian potential at the clock and $v$ is its total speed in that same frame \cite{IERS}. For a moving clock, write $\epsilon\equiv\delta\tau-1$ for the full inverse-rate excess; it has gravitational and kinematic entries at this order. Only in the static case does the exact lapse-gradient identity directly identify its spatial slope with the gravitational field. For a moving clock, differentiating the full rate ledger also involves the motion and the chosen trajectory; that full gradient must not be labeled gravitational acceleration without further analysis.

For an illustrative idealized Sun--Earth calculation, take a supported clock at Earth's surface, ignore Earth's rotation and the Sun's recoil, model Earth's orbit as circular at one astronomical unit, and set the potential zero at infinity. The solar potential, Earth's orbital motion, and Earth's surface potential contribute respectively
\begin{align}
  \epsilon_{\odot}
  &\simeq\frac{GM_{\odot}}{r_{\mathrm{AU}}c^2}
    =9.871\ \mathrm{ppb},\\
  \epsilon_{\mathrm{orb}}
  &\simeq\frac{v_{\mathrm{orb}}^2}{2c^2}
    =\frac{GM_{\odot}}{2r_{\mathrm{AU}}c^2}
    =4.935\ \mathrm{ppb},\\
  \epsilon_{\oplus}
  &\simeq\frac{GM_{\oplus}}{R_{\oplus}c^2}
    =0.696\ \mathrm{ppb},
  \label{eq:ledger-terms}
\end{align}
where $1\ \mathrm{ppb}=10^{-9}$. The circular-orbit relation $v_{\mathrm{orb}}^2=GM_{\odot}/r_{\mathrm{AU}}$ makes the orbital term half the solar-potential term. The potential terms alone sum to $\epsilon_{\odot}+\epsilon_{\oplus}=10.567\ \mathrm{ppb}$. Including the orbital term gives the total inverse-rate excess:
\begin{equation}
  \epsilon_{\mathrm{total}}\equiv\epsilon_{\odot}+\epsilon_{\mathrm{orb}}+\epsilon_{\oplus}
  =15.502\ \mathrm{ppb}.
  \label{eq:ledger-total}
\end{equation}
The parameters used are $GM_{\odot}=1.32712440041279419\times10^{20}\ \mathrm{m^3\,s^{-2}}$, $GM_{\oplus}=3.98600435507\times10^{14}\ \mathrm{m^3\,s^{-2}}$, $r_{\mathrm{AU}}=149\,597\,870\,700\ \mathrm{m}$, and the spherical mean-radius approximation $R_{\oplus}=6.371\times10^6\ \mathrm{m}$ \cite{Park,JPL}. The result is a model-dependent inverse-rate excess in the stated reference frame, not a reference-independent ``depth'' measured by the surface clock. To this order it corresponds to about $15.5\ \mathrm{ns}$ less proper time per reference second. This is a leading-order isolated Sun--Earth model, not an operational construction of an IAU coordinate time scale \cite{IERS}. Earth rotation, orbital eccentricity, other bodies, tides, barycentric corrections, and higher relativistic orders are excluded. The sum illustrates a clock-rate ledger; it is not an extra term in the static acceleration identity.

This example also shows why motion must be counted consistently. If a velocity is decomposed into multiple frame-dependent pieces, the cross terms in the squared total velocity cannot be discarded. One should evaluate the total $v^2$ once in a stated frame, alongside the total potential, rather than append independently named kinematic offsets without transforming the full model \cite{IERS}.

\section{Exact result in the Schwarzschild exterior}

For a spherical, nonrotating source of mass $M$, the exterior Schwarzschild line element is \cite{Schwarzschild,Wald,TaylorWheeler}
\begin{equation}
  \mathrm{d}s^2
  =-\left(1-\frac{r_s}{r}\right)c^2\mathrm{d}t^2
   +\left(1-\frac{r_s}{r}\right)^{-1}\mathrm{d}r^2
   +r^2\mathrm{d}\Omega^2,
  \qquad
  \mathrm{d}\Omega^2=\mathrm{d}\theta^2+\sin^2\theta\,\mathrm{d}\phi^2.
  \label{eq:schwarzschild-metric}
\end{equation}
The Schwarzschild radius and compactness are
\begin{equation}
  r_s\equiv\frac{2GM}{c^2},
  \qquad
  u\equiv\frac{GM}{rc^2}.
  \label{eq:schwarzschild-vars}
\end{equation}
In the vacuum exterior, the Schwarzschild metric has lapse
\begin{equation}
  \alpha(r)=\sqrt{1-\frac{r_s}{r}},
  \qquad
  \delta\tau(r)=\left(1-\frac{r_s}{r}\right)^{-1/2}.
  \label{eq:schwarzschild-rate}
\end{equation}
The proper radial distance at fixed time is $\mathrm{d}L=\mathrm{d}r/\alpha$. Thus Eq.~\eqref{eq:radial-gradient} gives the outward proper-radial gradient. Applying it to $\ln\delta\tau$ yields
\begin{equation}
  \nabla\ln\delta\tau
  =-\frac{GM}{c^2r^2\alpha(r)}\,\hat{\mathbf{r}},
  \qquad
  \boxed{\mathbf{g}
  =-\frac{GM}{r^2\alpha(r)}\,\hat{\mathbf{r}}.}
  \label{eq:schwarzschild-gradient}
\end{equation}
The proper acceleration required to hover has the opposite direction and magnitude
\begin{equation}
  a_{\mathrm{hold}}(r)
  =\frac{GM}{r^2\alpha(r)}
  =g_{\mathrm{N}}(r)\,\delta\tau(r),
  \qquad
  g_{\mathrm{N}}(r)\equiv\frac{GM}{r^2}.
  \label{eq:schwarzschild-hover}
\end{equation}
The reciprocal clock rate therefore multiplies the Newtonian acceleration in this exact Schwarzschild expression. This is a relation between a supported observer's local acceleration and the static lapse, not an assertion that a clock-rate ratio is a material density or an additional force.

Expanding in $u$ gives
\begin{equation}
  a_{\mathrm{hold}}
  =g_{\mathrm{N}}\left(
    1+u+\frac{3}{2}u^2+\frac{5}{2}u^3
      +\frac{35}{8}u^4+\cdots\right),
  \qquad |2u|<1.
  \label{eq:acceleration-series}
\end{equation}
This is the weak-field expansion of the exact acceleration magnitude, valid for $r>r_s$. The first correction to the Newtonian value is the compactness $u$. Near the horizon, the exact closed form is more informative than a truncated series.

Care is needed with radial coordinates. Direct differentiation with respect to the areal coordinate $r$ gives
\begin{equation}
  c^2\left|\frac{\mathrm{d}}{\mathrm{d}r}\ln\delta\tau\right|
  =\frac{g_{\mathrm{N}}}{\alpha^2},
  \label{eq:coordinate-gradient}
\end{equation}
which is not the local hovering acceleration. The proper spatial gradient includes the factor $\alpha$ in Eq.~\eqref{eq:radial-gradient}, giving $g_{\mathrm{N}}/\alpha$ instead. The distinction is the difference between change per unit coordinate radius and change per unit physical distance.

\section{Accelerometers, free fall, and the horizon}

A supported observer is not in free fall. A floor, rocket, or other support continually deflects the observer from a geodesic path, and an ideal accelerometer measures the proper acceleration $a_{\mathrm{hold}}$ in Eq.~\eqref{eq:schwarzschild-hover}. A freely falling observer follows a geodesic and reads zero proper acceleration, even at a location where the hovering acceleration would be large. Clock-rate gradients describe comparisons among static worldlines; an accelerometer reports the acceleration of its own worldline.

For $r>r_s$, the inverse rate and hovering acceleration are finite. As $r\to r_s^+$,
\begin{equation}
  \delta\tau\longrightarrow\infty,
  \qquad
  a_{\mathrm{hold}}\longrightarrow\infty.
  \label{eq:horizon-divergence}
\end{equation}
No static observer exists on the horizon. There is, however, a finite redshifted limit when the local hovering acceleration is multiplied by the lapse:
\begin{equation}
  \lim_{r\to r_s^+}\alpha(r)a_{\mathrm{hold}}(r)
  =\frac{GM}{r_s^2}
  =\frac{c^4}{4GM}.
  \label{eq:surface-gravity}
\end{equation}
With the time coordinate normalized at infinity, this is the Schwarzschild surface gravity expressed in acceleration units \cite{Wald,TaylorWheeler}. It is not the local acceleration of an observer at the horizon: the local hovering acceleration has no finite value there because a hovering observer cannot be placed there.

\section{Acceleration-associated scales}

An acceleration $g$ can be written as an inverse-time scale $g/c$ and a length scale $c^2/g$. For $g\simeq9.81\ \mathrm{m\,s^{-2}}$, these are approximately $32.7\ \mathrm{nHz}$ and $0.97$ light-years. The first is a dimensional frequency scale, not evidence that an Earth-bound observer emits or detects an oscillation at that frequency. The second is an acceleration length. For uniform proper acceleration in flat spacetime, the corresponding Rindler horizon lies a distance $c^2/g$ away; when $g$ varies in a gravitational field, this length is a local scale, not generally a global horizon distance \cite{Rindler}.

The Unruh temperature associated with uniform proper acceleration is
\begin{equation}
  T_{\mathrm{U}}=\frac{\hbar g}{2\pi c k_{\mathrm{B}}}.
  \label{eq:unruh-temperature}
\end{equation}
At terrestrial acceleration this is about $4\times10^{-20}\ \mathrm{K}$. Its energy expressed as an ordinary frequency is
\begin{equation}
  \frac{k_{\mathrm{B}}T_{\mathrm{U}}}{h}
  =\frac{g}{4\pi^2c},
  \label{eq:unruh-frequency}
\end{equation}
which is $4\pi^2$ times smaller than $g/c$. These equations are theoretical acceleration scales; they do not imply an observable clock signal at the inverse-rate gradient or a direct laboratory detection of an Earth-strength Unruh bath \cite{Unruh}.

\section{Conclusion}

For a static observer family, the spatial gradient of the inverse clock-rate ratio has a direct kinematic interpretation. With $\delta\tau=\mathrm{d}t/\mathrm{d}\tau$ and $\nabla$ defined using proper spatial distances, the inward gravitational direction is $\mathbf{g}=c^2\nabla\ln\delta\tau$ and the acceleration required to hold position is its negative. In the weak field this becomes Newtonian gravity. In the Schwarzschild exterior the hovering acceleration is exactly $GM/(r^2\alpha)=g_{\mathrm{N}}\delta\tau$.

The identity has a precise domain. It concerns static observers and a specified static spacetime; it is not a universal formula for arbitrary moving clocks. A moving-clock rate ledger includes both potential and kinetic contributions and depends on its reference frame. Clock comparisons can test the resulting relative rates, while accelerometers measure the proper acceleration of their own supported worldlines. Near a Schwarzschild horizon the hovering acceleration diverges and no static clock exists there, even though the appropriately redshifted surface-gravity limit is finite. In this formulation, ``temporal density'' remains a mnemonic for a rate ratio, and its gradient organizes established general-relativistic relationships without introducing new physics.

\clearpage
\begin{thebibliography}{99}
\small
\setstretch{1.0}
\raggedright
\urlstyle{same}

\bibitem{Schwarzschild}
K. Schwarzschild,
``\"Uber das Gravitationsfeld eines Massenpunktes nach der Einsteinschen Theorie,''
\emph{Sitzungsberichte der K\"oniglich Preussischen Akademie der Wissenschaften zu Berlin}, 189--196 (1916).

\bibitem{Wald}
R. M. Wald,
\emph{General Relativity}, chapter 12 (University of Chicago Press, Chicago, 1984).

\bibitem{TaylorWheeler}
E. F. Taylor and J. A. Wheeler,
\emph{Exploring Black Holes: Introduction to General Relativity}
(Addison Wesley Longman, San Francisco, 2000).

\bibitem{SantiagoVisser}
J. Santiago and M. Visser,
``Tolman-like temperature gradients in stationary spacetimes,''
\emph{Physical Review D} \textbf{98}, 064001 (2018),
\href{https://doi.org/10.1103/PhysRevD.98.064001}{doi:\nolinkurl{10.1103/PhysRevD.98.064001}}.

\bibitem{Chou}
C. W. Chou, D. B. Hume, T. Rosenband, and D. J. Wineland,
``Optical clocks and relativity,''
\emph{Science} \textbf{329}(5999), 1630--1633 (2010),
\href{https://doi.org/10.1126/science.1192720}{doi:\nolinkurl{10.1126/science.1192720}}.

\bibitem{Bothwell}
T. Bothwell et al.,
``Resolving the gravitational redshift across a millimetre-scale atomic sample,''
\emph{Nature} \textbf{602}, 420--424 (2022),
\href{https://doi.org/10.1038/s41586-021-04349-7}{doi:\nolinkurl{10.1038/s41586-021-04349-7}}.

\bibitem{IERS}
G. Petit and B. Luzum, eds.,
\emph{IERS Conventions (2010)}, IERS Technical Note 36, chapter 10 (2010),
\url{https://iers-conventions.obspm.fr/content/chapter10/tn36_c10.pdf}.

\bibitem{Park}
R. S. Park, W. M. Folkner, J. G. Williams, and D. H. Boggs,
``The JPL Planetary and Lunar Ephemerides DE440 and DE441,''
\emph{The Astronomical Journal} \textbf{161}, 105 (2021),
\href{https://doi.org/10.3847/1538-3881/abd414}{doi:\nolinkurl{10.3847/1538-3881/abd414}}.

\bibitem{JPL}
NASA/JPL Solar System Dynamics,
\emph{Astrodynamic Parameters}, \url{https://ssd.jpl.nasa.gov/astro_par.html}, and
\emph{Planetary Physical Parameters}, \url{https://ssd.jpl.nasa.gov/planets/phys_par.html},
accessed October 27, 2026.

\bibitem{Rindler}
W. Rindler,
``Kruskal space and the uniformly accelerated frame,''
\emph{American Journal of Physics} \textbf{34}, 1174--1178 (1966).

\bibitem{Unruh}
W. G. Unruh,
``Notes on black-hole evaporation,''
\emph{Physical Review D} \textbf{14}, 870--892 (1976),
\href{https://doi.org/10.1103/PhysRevD.14.870}{doi:\nolinkurl{10.1103/PhysRevD.14.870}}.

\end{thebibliography}

\end{document}